\section{Wasserstein (gradient) Flows}
\label{sec-wasserstein-flows}
\subsection{Evolutions over the Space of Measures}
\begin{equation}
\frac{\d x(t)}{\d t} = v_t(x(t)), \label{eq:lagrangian-advection}
\end{equation}
\begin{equation}
\frac{\partial \alpha_t}{\partial t} + \diverg(v_t \alpha_t) = 0. \label{eq:eulerian-advection}
\end{equation}
This PDE is only valid if $\alpha_t$ has a smooth density, and for general measures, this PDE \eqref{eq:eulerian-advection} should be understood in the weak sense, allowing it to be defined even for discrete measures with particles evolving according to \eqref{eq:lagrangian-advection}. Specifically, for any smooth test function $(t,x) \to \varphi(t,x)$ supported in time (in particular it should vanish at $t=0,1$),

\begin{equation}\label{eq:eulerian-advection-weak}
\int_0^1 \int_{\RR^d}
\left( \partial_t \varphi(t,x) + \langle v_t(x), \nabla_x \varphi(t,x) \rangle \right)
\d \alpha_t(x) \d t = 0.
\end{equation}
\begin{prop}[Lagrangian flows solve the continuity equation]
Consider a smooth flow $T_t: \RR^d \to \RR^d$ and denote $\alpha_t:= (T_t)_\sharp \alpha_0$. Define the Eulerian velocity field by \(v_t(T_t(y)):= \partial_t T_t(y).\)
Then $(\alpha_t, v_t)$ solves the continuity equation in the weak sense~\eqref{eq:eulerian-advection-weak}.
\end{prop}
\begin{proof}
Let $\varphi(t,x)$ be a smooth test function vanishing at $t=0$ and $t=1$. Using $\alpha_t=(T_t)_\sharp\alpha_0$ and the chain rule,
\[
\frac{\d}{\d t}\int \varphi(t,x)\d\alpha_t(x)
=
\frac{\d}{\d t}\int \varphi(t,T_t(y))\d\alpha_0(y)
=
\int \left(\partial_t \varphi(t,T_t(y))+
\dotp{\nabla_x\varphi(t,T_t(y))}{\partial_t T_t(y)}\right)\d\alpha_0(y).
\]
By the definition of $v_t$, this is
\[
\int \left(\partial_t \varphi(t,x)+
\dotp{\nabla_x\varphi(t,x)}{v_t(x)}\right)\d\alpha_t(x).
\]
Integrating in time and using the vanishing boundary values of $\varphi$ gives~\eqref{eq:eulerian-advection-weak}.
\end{proof}
\paragraph{From measure evolutions to vector fields.}
\begin{equation}
\partial_t \alpha_t + \diverg(v_t \alpha_t) = 0. \label{eq:inverse-flow}
\end{equation}
$$
\Hh_\alpha:= \{ v: \diverg( \alpha v ) = 0 \}.
$$
\paragraph{Dacorogna and Moser inversion.}
\begin{equation}\label{eq:dacorogna-moser}
v_t = - \frac{1}{\alpha_t} \nabla \Delta^{-1}(\partial_t \alpha_t).
\end{equation}
which was initially proposed by Dacorogna and Moser. A difficulty with this choice is that it is not well defined when $\alpha_t$ vanishes, and also it is not a gradient field, which might be desirable in some cases.

\paragraph{Least square inversion and gradient structure.}
\begin{equation}\label{eq:least-square-field}
\min_v \frac{1}{2}\int_0^1 \int_{\RR^d} \norm{v_t(x)}^2 \, \d\alpha_t(x) \, \d t \quad \text{subject to} \quad \diverg(\alpha v) + \frac{\partial \alpha}{\partial t} = 0.
\end{equation}
\begin{equation}
\min_v \max_{\psi} \frac{1}{2} \int_0^1 \left[ \int_{\RR^d} \norm{v_t(x)}^2 \, \d\alpha_t(x) + \int_{\RR^d} \psi_t(x)\left(\diverg(\alpha_t v_t)(x) + \partial_t \alpha_t(x)\right) \, \d x \right] \d t.
\end{equation}
\begin{equation}
\min_v \max_{\psi} \int_0^1 \left[ \int_{\RR^d} \left( \frac{1}{2} \norm{v_t(x)}^2 - \nabla \psi_t(x) \cdot v_t(x)\right) \, \d\alpha_t(x) + \int_{\RR^d} \psi_t(x) \, \partial_t \alpha_t(x) \, \d x \right] \d t.
\end{equation}
\(v_t(x) = \nabla \psi_t(x).\)
\begin{equation}\label{eq:least-square-field-explicit}
v_t = - \nabla \Delta_\alpha^{-1}(\partial_t \alpha_t), \quad \text{with} \quad \Delta_\alpha(\varphi):= \diverg(\alpha_t \nabla \varphi).
\end{equation}
\subsection{Generative Models via Flow Matching}
\paragraph{Stochastic interpolant.}
We assume that $\alpha_t$ is defined via a ``projection'' (in a loose sense) of a latent distribution $\pi \in \Pp(\RR^{d'})$, using an operator $P_t: \RR^{d'} \to \RR^d$ where $d' \gg d$, i.e.

\begin{equation}\label{eq:interp-coupling}
\forall t \in [0,1], \quad \alpha_t:= (P_t)_\sharp \pi.
\end{equation}
A common case is $d'=2d$. We write $(x,y) \in \RR^{d'} = \RR^d \times \RR^d$ and assume $P_0(x,y)=x$ and $P_1(x,y)=y$, so that $\pi$ is a probabilistic coupling between $\alpha_0$ and $\alpha_1$, i.e. $\pi$ has marginals $(\alpha_0, \alpha_1)$.

If $\pi = \alpha \otimes \beta$ and $\alpha = \frac{1}{n} \sum_i \delta_{x_i}$, $\beta = \frac{1}{m} \sum_j \delta_{y_j}$, then $\alpha_t$ consists of $n \times m$ Dirac masses

\(\alpha_t = \frac{1}{nm} \sum_{i,j} \delta_{P_t(x_i,y_j)}.\)
If $\pi = (\Id, T)_\sharp \alpha$ is a Brenier-type coupling, then $\alpha_t = ((1-t)\Id + tT)_\sharp \alpha$ is the so-called McCann OT interpolation.

\paragraph{Flow matching formula.}
\begin{prop}[Flow matching vector field]
For each fixed $t$, the solution of the following flow-matching problem
\begin{equation}
\min_{v_t} \int_{\RR^{d'}} \norm{v_t(P_t(u)) - [\partial_t P_t](u)}^2 \, \d\pi(u). \label{eq:flow-matching}
\end{equation}
or equivalently the conditional expectation
\begin{equation}\label{eq:flow-match-conditional}
v_t(z) = \EE_{u \sim \pi} \big( [\partial_t P_t](u) \, \big| \, z = P_t(u) \big).
\end{equation}
satisfies the continuity equation~\eqref{eq:inverse-flow}.
\end{prop}
\begin{proof} We now prove the flow matching formula \eqref{eq:flow-match-conditional}, i.e. that it defines a valid flow between $\alpha_0$ and $\alpha_1$.

First, let us recall the definition of the interpolated measure \(\alpha_t\) and the velocity field \(v_t\).

According to~\eqref{eq:interp-coupling}, the measure \(\alpha_t\) can be represented formally as: \(\alpha_t(z) = \int \delta(z - P_t(u)) \, \d\pi(u)\)
and rigorously for any test function \(\varphi(z)\) as:
\begin{equation}
\int \varphi(z) \, \d\alpha_t(z) = \int \varphi(P_t(u)) \, \d \pi(u). \label{eq:alpha_t}
\end{equation}
Following~\eqref{eq:flow-match-conditional}, the velocity field \(v_t\) can be written formally as: \(v_t(z) = \frac{1}{\alpha_t(z)} \int \delta(z - P_t(u)) [\partial_t P_t](u) \, \d\pi(u),\)
and rigorously for any vector field \(m(z)\) as:
\begin{equation}
\int \langle m(z), v_t(z) \rangle \, \d\alpha_t(z) = \int \langle m(P_t(u)), [\partial_t P_t](u) \rangle \, \d\pi(u). \label{eq:v_t}
\end{equation}
We aim to prove that the measure \(\alpha_t\) satisfies the continuity equation: \(\frac{\partial \alpha_t}{\partial t} + \diverg(\alpha_t v_t) = 0.\)
In a rigorous sense, this means showing that for any smooth test function \(\varphi(z)\),
\begin{equation}
\frac{\d}{\d t} \int \varphi(z) \, \d\alpha_t(z) - \int \langle v_t(z), \nabla \varphi(z) \rangle \, \d\alpha_t(z) = 0. \label{eq:continuity}
\end{equation}
To prove \eqref{eq:continuity}, we compute both terms separately and show that they cancel out. First, consider the time derivative of \(\int \varphi(z) \, \d\alpha_t(z)\). Using \eqref{eq:alpha_t}, we differentiate under the integral sign:
\(\frac{\d}{\d t} \int \varphi(z) \, \d\alpha_t(z) = \int \frac{\d}{\d t} \varphi(P_t(u)) \, \d\pi(u).\)
Applying the chain rule to \(t \to \varphi \circ P_t\):
\[
\frac{\d}{\d t} \varphi(P_t(u)) = \left\langle \nabla \varphi(P_t(u)), [\partial_t P_t](u) \right\rangle.
\]
Thus:
\begin{equation}
\frac{\d}{\d t} \int \varphi(z) \, \d\alpha_t(z) =
\int \left\langle \nabla \varphi(P_t(u)), [\partial_t P_t](u) \right\rangle \, \d\pi(u). \label{eq:time_derivative}
\end{equation}
Next, for the term involving \(\diverg(\alpha_t v_t)\), we use the definition of \(v_t\) in \eqref{eq:v_t} with \(m(z) = \nabla \varphi(z)\). Substituting this into the expression for \(v_t\), we get:
\(\int \langle v_t(z), \nabla \varphi(z) \rangle \, \d\alpha_t(z) = \int \langle \nabla \varphi(P_t(u)), [\partial_t P_t](u) \rangle \, \d \pi(u).\)
Comparing this result with \eqref{eq:time_derivative}, we see that: \(\frac{\d}{\d t} \int \varphi(z) \, \d\alpha_t(z) - \int \langle v_t(z), \nabla \varphi(z) \rangle \, \d\alpha_t(z) = 0.\)
\end{proof}
\paragraph{Connection with diffusion models.}
In the special case where $P_t(x,y)=(1-t)x+ty$ is a linear interpolation and $\pi = \alpha \otimes \beta$, then $\alpha_t$ is a convolution of rescaled versions of $\alpha_0$ and $\alpha_1$. Then the flow matching~\eqref{eq:flow-matching} boils down to

$$
\min_{(v_t)_t} \int_{\RR^{d} \times \RR^d} \norm{v_t( (1-t)x+t y ) - (y-x) }^2 \, \d\alpha_0(x) \d\alpha_1(y).
$$
If furthermore, $\alpha_0$ is an isotropic Gaussian, this is exactly (up to an exponential change of variable in the time variable) the diffusion model method.

\begin{proposition}[Tweedie identity]\label{prop:Tweedie}
Let $W$ be a random vector in $\RR^{d}$ with density $\beta$.
For $\sigma>0$, observe
\(Z \;=\; W + \sigma\,\varepsilon, \quad\text{where } \varepsilon \sim \Gaussian(0,I_{d}) \text{ is independent of } W.\)
Denote by \(\beta_\sigma \;=\; \beta * \Gaussian\bigl(0,\sigma^{2}I_{d}\bigr)\)
the density of $Z$.
Then
\(\EE\bigl[\,W \mid Z=z\bigr] \;=\; z \;+\;\sigma^{2}\,\nabla \log \beta_\sigma(z) \qquad\text{for all } z \in \RR^{d}.\)
\end{proposition}
\begin{proof}
Bayes' rule gives the conditional density
$
p_{W|Z}(w\mid z)
= \dfrac{\beta(w)\,\varphi_\sigma(z-w)}{\beta_\sigma(z)}
$
with $\varphi_\sigma$ the $\Gaussian(0,\sigma^{2}I_{d})$ density.
Hence
\(\EE[W\mid Z=z] = \frac{1}{\beta_\sigma(z)} \int_{\RR^{d}} w\, \beta(w)\,\varphi_\sigma(z-w)\,\d w.\)
Integrating by parts in $w$ and using
$
\nabla_z\varphi_\sigma(z-w)
= -\sigma^{-2}(z-w)\,\varphi_\sigma(z-w)
$
yields
\(\nabla_z\beta_\sigma(z) = \int \beta(w)\,\nabla_z\varphi_\sigma(z-w)\,\d w = -\sigma^{-2}\Bigl(z-\EE[W\mid Z=z]\Bigr)\,\beta_\sigma(z).\)
Rearranging finishes the proof.
\end{proof}
\begin{proposition}[Optimal vector field]\label{prop:flow}
Let $X\sim\alpha$ and $Y\sim\Gaussian(0,I_{d})$ be independent.
For $t\in(0,1)$ set
\(Z_t \;=\; (1-t)\,X + t\,Y, \qquad \alpha_t =\text{Law}(Z_t).\)
The minimizer $v^\star$ of
\(\min_{v}\;\int_{0}^{1}\! \iint_{\RR^{d}\times\RR^{d}} \bigl|y-x-v\bigl((1-t)x+t y,t\bigr)\bigr|^{2}\, \d\alpha(x)\,\d\Gaussian(y)\,\d t\)
is \(v^\star(x,t) = -\frac{1}{1-t}\,x \;-\; \frac{t}{1-t}\,\nabla\log\alpha_t(x) \qquad (x\in\RR^{d},\;t\in(0,1)).\)
\end{proposition}
\begin{proof}
Fix $t\in(0,1)$ and write $W=(1-t)X$, $\sigma=t$, so that
$Z_t = W + \sigma\,Y$ matches the setting of Proposition~\ref{prop:Tweedie}.
Conditional expectations satisfy
$
v^\star(z,t)
= \EE[Y-X\mid Z_t=z]
= \frac{1}{t}\,\EE[Z_t-W\mid Z_t=z]
-\,\frac{1}{1-t}\,\EE[W\mid Z_t=z].
$
Applying Proposition~\ref{prop:Tweedie} to $\EE[W\mid Z_t=z]$ and
noting $\EE[Y\mid Z_t=z]
= -\,t\,\nabla\log\alpha_t(z)$
gives the claimed formula.
\end{proof}
\paragraph{When is the induced map optimal?}
Integrating the learned velocity gives a deterministic map from $\alpha_0$ to $\alpha_1$, but this map is not automatically the Brenier optimal map. It is optimal only in special cases where the accumulated flow remains the gradient of a convex potential.

\subsection{Benamou--Brenier dynamic formulation of OT}
\begin{thm}[Benamou--Brenier]
For probability measures $\alpha_0,\alpha_1$ with finite second moments,
\begin{equation}\label{eq:benamou-brenier}
\Wass_2^2(\alpha_0,\alpha_1)
= \inf_{(\alpha_t,v_t)} \int_0^1 \int_{\RR^d} |v_t(x)|^2\, \d\alpha_t(x)\, \d t,
\end{equation}
where the infimum is over $(\alpha_t,v_t)$ solving
$\partial_t\alpha_t+\nabla\!\cdot(\alpha_t v_t)=0$ with $\alpha_{t=0}=\alpha_0$ and $\alpha_{t=1}=\alpha_1$.
If $\alpha_0$ has a density and $T$ is the optimal Monge map $T_{\sharp}\alpha_0=\alpha_1$, the minimizer is
\begin{equation}\label{eq:static-to-dynamic}
\alpha_t=((1-t)\Id+tT)_{\sharp}\alpha_0,
\qquad
v_t\bigl((1-t)x+tT(x)\bigr)=T(x)-x,\quad t\in[0,1].
\end{equation}
\end{thm}
\begin{proof}
\textit{Dynamic $\le$ static.} Let $T$ solve the Monge problem and define $(\alpha_t,v_t)$ by \eqref{eq:static-to-dynamic}. As $\dot T_t(x)=T(x)-x$ is $t$-independent,
\(\int_0^1\!\int |v_t|^2\,\d\alpha_t\,\d t=\int |T(x)-x|^2\,\d\alpha_0(x),\)
so the dynamic cost is no larger.

\smallskip
\textit{Static $\le$ dynamic.} Conversely, for a smooth deterministic path take the flow $T_t$ defined by $\dot T_t=v_t\!\circ\!T_t$, $T_0=\Id$. Then $\alpha_t=(T_t)_{\sharp}\alpha_0$ and $(T_1)_{\sharp}\alpha_0=\alpha_1$. Jensen's inequality gives
\(|T_1(x)-x|^2\le\int_0^1|v_t(T_t(x))|^2\,\d t.\)
Integrating yields
\(\int|T_1(x)-x|^2\,\d\alpha_0\le\int_0^1\!\int|v_t|^2\,\d\alpha_t\,\d t.\)
Thus the Monge cost is no larger than the dynamic one. For general finite-energy solutions of the continuity equation, the same argument is obtained by applying the superposition principle to lift the curve to a probability measure on absolutely continuous paths; this is the path-space formulation stated below.
\end{proof}
\begin{rem}[Path-space formulation]
Let $\Ss=C([0,1];\RR^d)$ be the space of continuous paths endowed with the uniform topology. For $t\in[0,1]$ define the evaluation map

\(P_t: \Ss\to\RR^d,\qquad P_t(\gamma)=\gamma(t).\)
\[
\Wass_2^2(\alpha_0,\alpha_1)=\inf_{M\in\Pp(\Ss)}\Bigl\{\int_{\Ss}\!\int_0^1|\dot\gamma(t)|^2\,\d t\,\d M(\gamma): (P_0)_{\sharp}M=\alpha_0,\ (P_1)_{\sharp}M=\alpha_1\Bigr\}.
\]
If $\alpha_0$ has a density, the minimizer $M^*$ is unique. Its time marginals reproduce the optimal curve: $\alpha_t=(P_t)_{\sharp}M^*$ for all $t$. Furthermore, denoting $Q_t: \Ss\to\RR^d, Q_t(\gamma)=\dot\gamma(t)$, the conditional law of the velocity is deterministic:

\((P_t, Q_t)_{\sharp}M^* = \alpha_t\otimes\delta_{v_t^*(x)},\qquad t\in[0,1],\)
where $v_t^*$ is the optimal velocity field in the Benamou--Brenier formulation. Hence $M^*$ concentrates on straight-line geodesics, assigning exactly one direction at each spatial point.

\end{rem}
\subsection{Wasserstein Gradient Flows}
\begin{equation}
\alpha_{t+\tau}:= \arg\min_\alpha \frac{1}{2 \tau} \Wass_2(\alpha_t, \alpha)^2 + f(\alpha). \label{eq:jko-discr}
\end{equation}
\paragraph{Euclidean gradient flows.}
If we restrict \eqref{eq:jko-discr} to finite dimensions and assume $\alpha_t = \delta_{x(t)}$ and $\alpha = \delta_x$ (single Dirac measures), this matches the implicit Euler scheme:

\(x(t+\tau):= \arg\min_x \frac{1}{2 \tau} \norm{x - x(t)}^2 + h(x),\)
where $h(x) = f(\delta_x)$. Its solution is formally given by the implicit Euler formula:

\(x(t+\tau) = (\Id + \tau \nabla h)^{-1}(x(t)).\)
\(x(t+\tau) = (\Id - \tau \nabla h)(x(t)) = x(t) - \tau \nabla h(x(t)).\)
\begin{equation}
\dot{x}(t) = -\nabla h(x(t)). \label{eq:grad-flow-classical}
\end{equation}
\paragraph{Wasserstein gradient formula.}
As $\tau \to 0$, under certain conditions on $f$, \eqref{eq:jko-discr} defines a continuous evolution $t \mapsto \alpha_t$. As discussed earlier, this evolution can be described as a Lagrangian evolution \eqref{eq:lagrangian-advection}.

\begin{equation}
\frac{\partial \alpha_t}{\partial t} + \diverg(-\Wgrad f(\alpha_t) \alpha_t) = 0. \label{eq:wassflow-pde}
\end{equation}
\(\Wgrad f(\alpha) = \nabla_x \varphi, \quad \text{where } \varphi:= \delta f(\alpha).\)
\(f((1-\tau)\alpha + \tau \beta) = f(\alpha + \tau \rho) = f(\alpha) + \tau \int [\delta f(\alpha)](x) \d \rho(x) + o(\tau),\)
where $\rho = \beta - \alpha$ is a zero-mean measure.

\begin{proposition}[Formal Wasserstein gradient]\label{prop-formal-wass-gradient}
Assume that $f$ admits a smooth first variation $\delta f(\alpha)$ and that $\alpha$ has a smooth positive density. For infinitesimal perturbations generated by a velocity field $v$ through $(\Id+\tau v)_\sharp\alpha$, the differential of $f$ is
\(\frac{\d}{\d\tau}_{|\tau=0} f\big((\Id+\tau v)_\sharp\alpha\big) = \int \dotp{\nabla \delta f(\alpha)(x)}{v(x)}\d\alpha(x).\)
Hence, for the Riemannian metric $\norm{v}_{L^2(\alpha)}^2=\int \norm{v}^2\d\alpha$, the Wasserstein gradient is the vector field
\(\Wgrad f(\alpha)=\nabla \delta f(\alpha).\)
\end{proposition}
\begin{proof}
The push-forward expansion gives, in the sense of distributions, \((\Id+\tau v)_\sharp\alpha = \alpha-\tau\diverg(\alpha v)+o(\tau).\)
Using the definition of the first variation, \(f((\Id+\tau v)_\sharp\alpha) = f(\alpha)-\tau\int \delta f(\alpha)\diverg(\alpha v)\d x+o(\tau).\)
An integration by parts, with either compact support or vanishing boundary flux, gives \(-\int \delta f(\alpha)\diverg(\alpha v)\d x = \int \dotp{\nabla\delta f(\alpha)}{v}\d\alpha.\)
By definition of the Riesz representative for the $L^2(\alpha)$ metric, this representative is $\nabla\delta f(\alpha)$.
\end{proof}
\paragraph{From the JKO step to the velocity field.}
$$
\min_{v} \frac{1}{2 \tau} \tau^2 \norm{v}_{L^2(\alpha_t)}^2 + f( (\Id + \tau v)_\sharp \alpha_t ).
$$
$$
(\Id + \tau v)_\sharp \alpha_t = \alpha_t - \tau \diverg( v \alpha_t ) + o(\tau)
$$
$$
f( (\Id + \tau v)_\sharp \alpha_t ) = f(\alpha_t) - \tau \int \delta f(\alpha_t) \diverg( v \alpha_t ) \d x + o(\tau)
$$
$$
= f(\alpha_t) + \tau \int \langle \nabla_x \delta f(\alpha_t)(x), v(x) \rangle \d \alpha_t(x) + o(\tau)
$$
$$
\min_{v} f(\alpha_t) + \tau \int \Big[ \frac{1}{2} \norm{v(x)}^2 + \langle \Wgrad f(\alpha_t)(x), v(x) \rangle \Big]\d \alpha_t(x) + o(\tau).
$$
\paragraph{Discrete evolutions.}
If $f(\alpha)$ can be evaluated on discrete distributions and $\Wgrad$ is continuous in this case, the flow \eqref{eq:wassflow-pde} maintains the number of Dirac masses, $\alpha_t = \frac{1}{n} \sum_i \delta_{x_i(t)}$. The particles $X(t):= (x_i(t))_i$ evolve according to a system of coupled ODEs:

\begin{equation}
\dot{x}_i(t) = -n\nabla_{x_i} F(X(t)), \label{eq:wassflows-particles}
\end{equation}
where $F(X):= f\left(\frac{1}{n} \sum_i \delta_{x_i}\right)$ and the factor $n$ comes from the empirical Wasserstein metric $\frac1n\sum_i \norm{\dot x_i}^2$.

\paragraph{Linear Functionals.}
\begin{equation}\label{eq:linear-func}
f(\alpha) = \int h(x) \d \alpha(x).
\end{equation}
\(\frac{\partial \alpha_t}{\partial t} + \diverg(-\nabla h \alpha_t) = 0.\)
\paragraph{Shannon Neg-Entropy.}
\begin{equation}
f(\alpha) = \int \log\left(\frac{\d \alpha}{\d x}(x)\right) \d \alpha(x). \label{eq:entropy-func}
\end{equation}
\(\partial_t \alpha_t = \Delta \alpha_t.\)
\begin{equation}\label{eq:gen-entropies}
f(\alpha) = \int g\left(\frac{\d \alpha}{\d x}\right) \d x,
\end{equation}
\(\frac{\partial \alpha_t}{\partial t} = \Delta(P(\alpha_t)),\)
where the pressure $P$ satisfies $P'(s)=s g''(s)$.

\paragraph{Interaction Energies.}
\begin{equation}
f(\alpha):= \iint k(x, y) \d \alpha(x) \d \alpha(y). \label{eq:quadratic-func}
\end{equation}
For a symmetric kernel $k$:

\(\delta f(\alpha)(x) = 2 \int k(x, y) \d \alpha(y), \quad \Wgrad f(\alpha)(x) = 2 \int \nabla_x k(x, y) \d \alpha(y).\)
For $\alpha_0 = \frac{1}{n} \sum_i \delta_{x_i}$, the flow \eqref{eq:wassflow-pde} implies particles $(x_i(t))_i$ obey:

\(\dot{x}_i(t) = -\frac{2}{n} \sum_j \nabla k(x_i(t), x_j(t)).\)
If $k$ is positive definite and one minimizes the squared MMD distance to a teacher distribution $\beta$, then

\[
\norm{\alpha-\beta}_k^2
=
\iint k\d\alpha\d\alpha
-2\int\!\left(\int k(x,y)\d\beta(y)\right)\d\alpha(x)
+\mathrm{constant}.
\]
Thus MMD training energies are exactly an interaction energy plus a linear potential; the teacher distribution appears through the potential $x\mapsto-2\int k(x,y)\d\beta(y)$.

\paragraph{Geodesics and geodesic convexity.}
\label{sec-geodesic-convexity}
\(\alpha_t=((1-t)P_0+tP_1)_\sharp\pi^\star, \qquad t\in[0,1],\)
where $P_0(x,y)=x$ and $P_1(x,y)=y$. If the optimal plan is induced by a Brenier map $T$, this reduces to $((1-t)\Id+tT)_\sharp\alpha_0$.

\begin{defn}[Geodesic convexity]\label{def-geodesic-convexity}
A functional $f$ on $\Pp_2(\RR^d)$ is geodesically convex if for every $\Wass_2$ geodesic $(\alpha_t)_t$, \(f(\alpha_t)\leq (1-t)f(\alpha_0)+t f(\alpha_1).\)
It is $\lambda$-geodesically convex if the right-hand side is improved by $-\frac{\lambda}{2}t(1-t)\Wass_2^2(\alpha_0,\alpha_1)$.
\end{defn}
\begin{prop}[Basic geodesically convex energies]\label{prop-basic-geodesic-convexity}
\begin{enumerate}
\item If $h$ is convex, then $\alpha\mapsto\int h\d\alpha$ is geodesically convex; if $h$ is $\lambda$-strongly convex, it is $\lambda$-geodesically convex.
\item If $W(x-y)$ is convex as a function of the displacement, then $\alpha\mapsto\frac12\iint W(x-y)\d\alpha(x)\d\alpha(y)$ is geodesically convex.
\item Shannon entropy $\alpha\mapsto\int\rho\log\rho\,\d x$ is geodesically convex.
\item The relative entropy $\KL(\alpha|\gamma)$ with $\d\gamma=e^{-V}\d x/Z$ is $\lambda$-geodesically convex when $V$ is $\lambda$-strongly convex.
\end{enumerate}
\end{prop}
\begin{proof}
Along a Monge geodesic $X_t=(1-t)X_0+tX_1$, convexity of $h$ gives $h(X_t)\leq(1-t)h(X_0)+t h(X_1)$, and strong convexity gives the additional quadratic term; integrating proves the first claim. The interaction claim follows similarly by applying convexity of $W$ to pairwise differences $X_t-X_t'=(1-t)(X_0-X_0')+t(X_1-X_1')$ and integrating over two independent copies. The entropy claim is McCann's displacement convexity theorem; at the density level it follows from the concavity of the Jacobian determinant under the interpolation of optimal maps. Finally, $\KL(\alpha|\gamma)=\int\rho\log\rho\,\d x+\int V\d\alpha+\mathrm{constant}$, so it is the sum of displacement-convex entropy and a $\lambda$-geodesically convex linear potential.
\end{proof}
\paragraph{Convergence of the flow.}
\begin{prop}[Energy decay for convex Wasserstein flows]\label{prop-convex-wass-flow-rate}
Assume formally that $f$ is geodesically convex, admits a smooth first variation, and has a minimizer $\alpha^\star$. Let $(\alpha_t)_t$ be a smooth solution of the Wasserstein gradient flow
\(\partial_t\alpha_t+\diverg(\alpha_t v_t)=0, \qquad v_t=-\Wgrad f(\alpha_t).\)
Then
\(\frac{\d}{\d t} f(\alpha_t) = -\int \norm{\Wgrad f(\alpha_t)(x)}^2\,\d\alpha_t(x) \leq 0.\)
If $T_t$ is the optimal map from $\alpha_t$ to $\alpha^\star$, then
\(f(\alpha_t)-f(\alpha^\star) \leq -\frac{\d}{\d t}\frac12 \Wass_2^2(\alpha_t,\alpha^\star),\)
and consequently
\(f(\alpha_t)-f(\alpha^\star) \leq \frac{\Wass_2^2(\alpha_0,\alpha^\star)}{2t}.\)
If $f$ is $\lambda$-geodesically convex with $\lambda>0$, then
\(f(\alpha_t)-f(\alpha^\star) \leq e^{-2\lambda t}\bigl(f(\alpha_0)-f(\alpha^\star)\bigr).\)
\end{prop}
\begin{proof}
The chain rule and Proposition~\ref{prop-formal-wass-gradient} give
\(\frac{\d}{\d t}f(\alpha_t) = \int \dotp{\Wgrad f(\alpha_t)(x)}{v_t(x)}\,\d\alpha_t(x) = -\int \norm{\Wgrad f(\alpha_t)(x)}^2\,\d\alpha_t(x).\)
Geodesic convexity along the geodesic
$((1-s)\Id+sT_t)_\sharp\alpha_t$ gives \(f(\alpha^\star)-f(\alpha_t) \geq \int \dotp{\Wgrad f(\alpha_t)(x)}{T_t(x)-x}\,\d\alpha_t(x).\)
Since $v_t=-\Wgrad f(\alpha_t)$, this reads \(f(\alpha_t)-f(\alpha^\star) \leq \int \dotp{v_t(x)}{T_t(x)-x}\,\d\alpha_t(x).\)
\(\frac{\d}{\d t}\frac12\Wass_2^2(\alpha_t,\alpha^\star) = \int \dotp{x-T_t(x)}{v_t(x)}\,\d\alpha_t(x),\)
which proves the differential inequality. Integrating it from $0$ to $t$ and using the monotonicity of $s\mapsto f(\alpha_s)$ gives
\[
t\bigl(f(\alpha_t)-f(\alpha^\star)\bigr)
\leq
\int_0^t \bigl(f(\alpha_s)-f(\alpha^\star)\bigr)\,\d s
\leq
\frac12\Wass_2^2(\alpha_0,\alpha^\star).
\]
If $f$ is $\lambda$-geodesically convex, the Wasserstein analogue of strong convexity gives the slope inequality
\(\int\norm{\Wgrad f(\alpha_t)}^2\d\alpha_t \geq 2\lambda\bigl(f(\alpha_t)-f(\alpha^\star)\bigr).\)
\(\frac{\d}{\d t}\bigl(f(\alpha_t)-f(\alpha^\star)\bigr) \leq -2\lambda\bigl(f(\alpha_t)-f(\alpha^\star)\bigr),\)
and Gronwall's lemma gives the exponential rate.
\end{proof}
\begin{prop}[Convex examples covered by the theory]\label{prop-convex-flow-examples}
The energy decay proposition applies to linear potentials~\eqref{eq:linear-func} when $h$ is convex, to interaction energies~\eqref{eq:quadratic-func} when the interaction potential is convex in displacements, and to entropy-type functionals such as Shannon entropy~\eqref{eq:entropy-func}.
\end{prop}
\begin{proof}
linear potentials inherit convexity from the pointwise potential, interaction energies inherit it from pairwise displacement convexity, and Shannon entropy is displacement convex by McCann's theorem.
\end{proof}
\subsection{Application: Training Two-Layer MLPs as Wasserstein Flows}
\(g_\theta(u):= \frac{1}{n} \sum_i \psi(\theta_i, u), \quad \text{where } \psi(\theta_i, u):= a_i \langle u, w_i \rangle,\)
and $\theta_i:= (w_i, a_i) \in \RR^{d+1}$ represents the parameters of a neuron. Importantly, these functions are invariant under permutations of the neurons.

\(G_\alpha(u):= \int_{\RR^{d+1}} \psi(\theta, u) \, \d \alpha(\theta).\)
\(\min_\alpha f(\alpha):= \frac{1}{N} \sum_{k=1}^N \ell(G_\alpha(u_k), y_k).\)
\begin{equation*}
\dot{\theta}_i = -n\nabla_{\theta_i} F(\theta), \quad \text{where } F(\theta):= f\left(\frac{1}{n} \sum_i \delta_{\theta_i}\right).
\end{equation*}
This is equivalent to the PDE \eqref{eq:wassflow-pde} on the neuron density, where the Wasserstein gradient is used.

\begin{equation*}
\delta f(\alpha)(\theta) = \frac{1}{N} \sum_{k=1}^N \left(\int \psi(\theta', u_k) \, \d \alpha(\theta') - y_k\right) \psi(\theta, u_k),
\end{equation*}
\(\delta f(\alpha)(\theta) = \int k(\theta, \theta') \, \d \alpha(\theta') + g(\theta),\)
\begin{align}
k(\theta, \theta') &:= \frac{1}{N} \sum_k \psi(\theta, u_k) \psi(\theta', u_k), \\
g(\theta) &:= -\frac{1}{N} \sum_k y_k \psi(\theta, u_k).
\end{align}
\(\Wgrad f(\alpha)(\theta) = \int \nabla_\theta k(\theta, \theta') \, \d \alpha(\theta') + \nabla_\theta g(\theta).\)
\subsection{Application: One-Step Generative Models}
Let $\zeta$ be a simple latent distribution and let $\alpha_\theta=(G_\theta)_\sharp\zeta$ be the model distribution. Assume that the target data distribution is $\beta$. A Wasserstein-flow construction chooses a discrepancy

\(\mathcal E_\beta(\alpha),\)
\begin{equation}\label{eq-one-step-wgf}
\partial_t\mu_t+\diverg(\mu_t w_t)=0,
\qquad
w_t(x)=-\nabla\delta_\alpha \mathcal E_\beta(\mu_t)(x).
\end{equation}
\begin{equation}\label{eq-one-step-l2-fit}
\min_\eta \int_0^1\!\int
\norm{U_\eta(t,x)-w_t(x)}^2
\,\d\mu_t(x)\,\d t.
\end{equation}
\[
\alpha_{\theta}^{+}
=
(\Id+\tau U_\eta)_\sharp \alpha_\theta,
\qquad\text{or equivalently}\qquad
G_\theta^{+}(z)=G_\theta(z)+\tau U_\eta(G_\theta(z)).
\]
\paragraph{Normalized drifting fields.}
\begin{equation}\label{eq-normalized-kernel-drift}
B_\epsilon[\nu](x)
\eqdef
\frac{\int (y-x)K_\epsilon(x,y)\,\d\nu(y)}
{\int K_\epsilon(x,y)\,\d\nu(y)}.
\end{equation}
For the Gaussian kernel $K_\epsilon(x,y)=\exp(-\norm{x-y}^2/(2\epsilon))$, this normalized field is a score of a smoothed density:

\begin{equation}\label{eq-normalized-kernel-score}
B_\epsilon[\nu](x)
=
\epsilon\nabla\log\!\left(\int K_\epsilon(x,y)\,\d\nu(y)\right).
\end{equation}
\begin{equation}\label{eq-cross-minus-self-drift}
u_t(x)=B_\epsilon[\beta](x)-B_\epsilon[\mu_t](x)
=
\epsilon\nabla\log
\frac{\int K_\epsilon(x,y)\,\d\beta(y)}
{\int K_\epsilon(x,y)\,\d\mu_t(y)}.
\end{equation}
\begin{prop}[Drifting as a time-dependent Wasserstein gradient]\label{prop-drifting-semi-relaxed-gradient}
Let $\mu_t$ be a smooth curve of positive densities and let $u_t=\nabla\phi_t$ be a smooth time-dependent gradient field. Define the semi-relaxed functional
\begin{equation}\label{eq-semi-relaxed-drift-functional}
\mathcal R_t(\alpha|\mu_t)
\eqdef
-\int \phi_t(x)\,\d\alpha(x)
+\int \phi_t(x)\,\d\mu_t(x).
\end{equation}
Here $\mu_t$ and $\phi_t$ are frozen when taking the first variation with respect to the first argument $\alpha$. Then the continuity equation
\(\partial_t\mu_t+\diverg(\mu_t u_t)=0\)
is the formal Wasserstein gradient descent of the time-dependent functional $\alpha\mapsto\mathcal R_t(\alpha|\mu_t)$.
\end{prop}
\begin{proof}
Since $\mu_t$ and $\phi_t$ are fixed in the variation with respect to $\alpha$, the first variation is \(\delta_\alpha \mathcal R_t(\alpha|\mu_t)(x)=-\phi_t(x).\)
By Proposition~\ref{prop-formal-wass-gradient}, \(\Wgrad \mathcal R_t(\alpha|\mu_t) = \nabla\delta_\alpha \mathcal R_t(\alpha|\mu_t) = -\nabla\phi_t = -u_t.\)
The Wasserstein gradient-descent velocity is the negative of this gradient, namely $u_t$. Substituting this velocity in the continuity equation gives the claimed flow.
\end{proof}
\begin{example}[Kernel drifting as a semi-relaxed divergence]
For the Gaussian-kernel drift~\eqref{eq-cross-minus-self-drift}, set

\(\phi_t(x)= \epsilon\log \frac{\int K_\epsilon(x,y)\,\d\beta(y)} {\int K_\epsilon(x,y)\,\d\mu_t(y)}.\)
Then $u_t=\nabla\phi_t$, so Proposition~\ref{prop-drifting-semi-relaxed-gradient} shows that kernel drifting is the Wasserstein gradient descent of

\[
\mathcal R_t^{\mathrm{drift}}(\alpha|\mu_t)
=
\epsilon
\int
\log
\frac{\int K_\epsilon(x,y)\,\d\mu_t(y)}
{\int K_\epsilon(x,y)\,\d\beta(y)}
\,\d\alpha(x)
+\mathrm{constant}.
\]
\end{example}
\begin{rem}[General fields and projection onto gradients]
\(\nabla\phi_t = \uargmin{\nabla\phi} \int \norm{\nabla\phi(x)-b_t(x)}^2\,\d\mu_t(x).\)
\end{rem}
\subsection{Application: Evolution in Depth of Transformers}
\paragraph{Attention as a context-dependent velocity.}
After tokenization, embedding, and positional encoding, each input (from a set of tokens) is represented as a point cloud $(x_i)_{i=1}^n$ of $n$ points in the space of vectorized tokens. An attention layer with skip connection and rescaling by $1/T$ (where $T$ is the depth) defines a transformation of the tokens:

\(x_i \mapsto x_i + \frac{1}{T} \sum_j \frac{e^{\langle Q x_i, K x_j \rangle} V x_j}{\sum_{\ell} e^{\langle Q x_i, K x_\ell \rangle}},\)
where $\theta = (K, Q, V)$ are the parameters of the attention layer, represented by three matrices.

\paragraph{Token measure evolution.}
To handle an arbitrary number of tokens, we define $\alpha = \frac{1}{n} \sum_i \delta_{x_i}$ as the empirical measure of tokens and rewrite the transformer mapping as:

\(x_i \mapsto x_i + \frac{1}{T} \Gamma_\theta[\alpha](x_i),\)
\(\Gamma_\theta[\alpha](x):= \int \frac{e^{\langle Q x, K y \rangle} V y \, \d \alpha(y)}{\int e^{\langle Q x, K z \rangle} \, \d \alpha(z)}.\)
\(\alpha_{t+\tau} = (\Id + \tau \Gamma_{\theta_t}[\alpha_t])_\sharp \alpha_t.\)
\(\partial_t \alpha_t + \diverg(\alpha_t \Gamma[\alpha_t]) = 0.\)
\paragraph{Gradient structure and limitations.}
When $V=Q^\top K$, the field $\Gamma[\alpha]$ is a gradient vector field, since it is the gradient of the log-partition function associated with the attention weights. It is not, however, the gradient of a first variation, so this PDE is not a Wasserstein gradient flow.

\subsection{Gaussian Closures as Finite-Dimensional Test Cases}
\begin{prop}[Affine velocities preserve Gaussianity]\label{prop-gaussian-affine-closure}
Let $\alpha_t=\Gaussian(\mean_t,\cov_t)$, with $\cov_t$ positive definite, solve the continuity equation with an affine velocity
\(v_t(x)=b_t+A_t(x-\mean_t).\)
Then $\alpha_t$ remains Gaussian and its moments solve
\(\dot \mean_t=b_t, \qquad \dot\cov_t=A_t\cov_t+\cov_t A_t^\top.\)
Conversely, any smooth Gaussian curve with positive definite covariance can be generated by such an affine velocity. If one wants the velocity to be a Wasserstein tangent gradient, one chooses the unique symmetric solution of the Lyapunov equation
\(A_t\cov_t+\cov_t A_t=\dot\cov_t.\)
\end{prop}
\begin{proof}
Let $X_t$ follow the characteristic ODE $\dot X_t=b_t+A_t(X_t-\mean_t)$. This linear ODE maps Gaussian random variables to Gaussian random variables. Taking expectation gives $\dot\mean_t=b_t$. Writing $\tilde X_t=X_t-\mean_t$, one has $\dot{\tilde X}_t=A_t\tilde X_t$, hence
\(\dot\cov_t = \frac{\d}{\d t}\EE(\tilde X_t\tilde X_t^\top) = A_t\cov_t+\cov_t A_t^\top.\)
For the converse, set $b_t=\dot\mean_t$ and choose any matrix $A_t$ satisfying the covariance equation. Since $\cov_t$ is positive definite, the Lyapunov map $A\mapsto A\cov_t+\cov_t A$ is invertible on symmetric matrices, which gives the unique symmetric choice when a gradient velocity is required. In that case $v_t$ is the gradient of the quadratic potential $x\mapsto \dotp{b_t}{x}+\dotp{A_t(x-\mean_t)}{x-\mean_t}/2$.
\end{proof}
\begin{example}[Flow matching and diffusion paths between Gaussians]
\(v_t(x)=\dot\mean_t+A_t(x-\mean_t), \qquad A_t\cov_t+\cov_t A_t=\dot\cov_t.\)
\(v_t(x)=\dot\mean_t+\frac{\dot s_t}{s_t}(x-\mean_t).\)
\(X_t=a_tX_0+\sigma_t Z,\qquad Z\sim\Gaussian(0,\Id),\)
has $\mean_t=a_t\mean_0$ and $\cov_t=a_t^2\cov_0+\sigma_t^2\Id$. Thus, in the Gaussian case, diffusion paths and flow-matching paths reduce to the same mean-covariance bookkeeping, although the corresponding training objectives are different.

\end{example}
\begin{rem}[Gaussian Wasserstein-gradient closures]
Let $F(\mean,\cov)$ be the restriction of a smooth functional $f$ to Gaussian measures. The Wasserstein gradient flow constrained to the Gaussian family is formally

\begin{equation}\label{eq-gaussian-wgf-closure}
\dot\mean_t=-\nabla_\mean F(\mean_t,\cov_t),
\qquad
\dot\cov_t=-2\bigl(\cov_t\nabla_\cov F(\mean_t,\cov_t)+\nabla_\cov F(\mean_t,\cov_t)\cov_t\bigr),
\end{equation}
where $\nabla_\cov F$ is the symmetric matrix derivative. For the KL flow toward $\beta=\Gaussian(\bar\mean,\bar\cov)$, this gives the Ornstein--Uhlenbeck equations

\(\dot\mean_t=-\bar\cov^{-1}(\mean_t-\bar\mean), \qquad \dot\cov_t=2\Id-\bar\cov^{-1}\cov_t-\cov_t\bar\cov^{-1}.\)
For the Sinkhorn divergence between Gaussian measures, one can similarly restrict the closed Gaussian formula to the variables $(\mean,\cov)$.

\end{rem}
\begin{rem}[Other Gaussian test energies]
\end{rem}
\begin{example}[Linear mean-field networks]
\(S=\frac1N\sum_{k=1}^N u_ku_k^\top, \qquad r=\frac1N\sum_{k=1}^N y_ku_k, \qquad q_t=\int aw\,\d\alpha_t(a,w).\)
For the square loss, set $h_t=S q_t-r$. The first variation is

\(\delta f(\alpha_t)(a,w)=a\dotp{w}{h_t}.\)
\[
-\nabla_{(a,w)}\delta f(\alpha_t)(a,w)
=
-\begin{pmatrix}
0 & h_t^\top \\
h_t & 0
\end{pmatrix}
\begin{pmatrix} a \\ w \end{pmatrix}.
\]
\end{example}
\begin{example}[Gaussian kernel drifting]
Let the target be $\beta=\Gaussian(\bar\mean,\bar\cov)$ and assume $\mu_t=\Gaussian(\mean_t,\cov_t)$. For the Gaussian kernel

\(K_\epsilon(x,y)=\exp(-\norm{x-y}^2/(2\epsilon)),\)
the normalized field~\eqref{eq-normalized-kernel-drift} satisfies

\(B_\epsilon[\mu_t](x) = -\epsilon(\cov_t+\epsilon\Id)^{-1}(x-\mean_t).\)
Thus the drifting velocity~\eqref{eq-cross-minus-self-drift} is affine and preserves Gaussianity. With

\(A_t=(\cov_t+\epsilon\Id)^{-1}, \qquad \bar A=(\bar\cov+\epsilon\Id)^{-1},\)
\(\dot\mean_t=\epsilon\bar A(\bar\mean-\mean_t), \qquad \dot\cov_t=\epsilon\bigl((A_t-\bar A)\cov_t+\cov_t(A_t-\bar A)\bigr).\)
\end{example}
\begin{example}[Gaussian closure of attention dynamics]
For the transformer PDE, assume $\alpha=\Gaussian(\mean,\cov)$. Since exponential tilting preserves Gaussianity,

\(\frac{\int e^{\dotp{Qx}{Ky}}\,y\,\d\alpha(y)} {\int e^{\dotp{Qx}{Kz}}\,\d\alpha(z)} = \mean+\cov K^\top Qx.\)
\(\Gamma_\theta[\alpha](x)=V\mean+V\cov K^\top Qx\)
\(\dot\mean_t=(V_t+V_t\cov_tK_t^\top Q_t)\mean_t, \qquad \dot\cov_t=B_t\cov_t+\cov_tB_t^\top, \qquad B_t=V_t\cov_tK_t^\top Q_t.\)
When $V_t=Q_t^\top K_t$, the matrix $B_t=Q_t^\top K_t\cov_tK_t^\top Q_t$ is symmetric positive semidefinite, matching the gradient-field case mentioned above. It is instead a tractable model of how attention can shear, amplify or contract a cloud of tokens through its covariance.

\end{example}